\documentclass{article}
\usepackage{graphicx}
\usepackage{geometry}
\geometry{a4paper, margin=1in}

\usepackage{listings}
\usepackage{xcolor}
\usepackage{tcolorbox}

\definecolor{codeGreen}{rgb}{0,0.6,0}
\definecolor{codeGray}{rgb}{0.5,0.5,0.5}
\definecolor{codePurple}{rgb}{0.58,0,0.82}
\definecolor{codeBg}{rgb}{0.95,0.95,0.92}

\definecolor{terminalBg}{rgb}{0.15,0.15,0.15}
\definecolor{terminalText}{rgb}{0.95,0.95,0.95}

\lstdefinestyle{pythonStyle}{
    backgroundcolor=\color{codeBg},   
    commentstyle=\color{codeGreen},
    keywordstyle=\color{magenta},
    numberstyle=\tiny\color{codeGray},
    stringstyle=\color{codePurple},
    basicstyle=\ttfamily\footnotesize,
    breakatwhitespace=false,         
    breaklines=true,                 
    captionpos=b,                    
    keepspaces=true,                 
    numbers=left,                    
    numbersep=5pt,                  
    showspaces=false,                
    showstringspaces=false,
    showtabs=false,                  
    tabsize=4
}

\lstdefinestyle{bashStyle}{
    backgroundcolor=\color{terminalBg},   
    basicstyle=\ttfamily\footnotesize\color{terminalText},
    keywordstyle=\color{green},
    stringstyle=\color{cyan},
    breakatwhitespace=false,         
    breaklines=true,                 
    keepspaces=true,                 
    numbers=none,
    showspaces=false,                
    showstringspaces=false,
    showtabs=false,                  
    tabsize=2
}

\lstdefinestyle{xmlStyle}{
    backgroundcolor=\color{codeBg},   
    basicstyle=\ttfamily\footnotesize\color{black},
    moredelim=[s][\color{blue}]{<}{>},
    moredelim=[s][\color{codeGreen}]{"}{"},
    showspaces=false,                
    showstringspaces=false,
    showtabs=false,
    breaklines=true,                 
    keepspaces=true,                 
    numbers=none,                  
    tabsize=2
}

\newcommand{\code}[1]{\colorbox{codeBg}{\texttt{#1}}}

\newtcolorbox{notebox}[1][Note]{
    colback=blue!5!white,
    colframe=blue!80!black,
    fonttitle=\bfseries,
    title=#1,
    arc=2mm
}

\newtcolorbox{warnbox}[1][Warning]{
    colback=red!5!white,
    colframe=red!75!black,
    fonttitle=\bfseries,
    title=#1,
    arc=2mm
}

\title{ROS2 Tutorials - ROS2 Humble For Beginners}
\author{Piotr Sarna}
\date{March 2026}

\begin{document}

\maketitle
\tableofcontents
\newpage

\section{Intro: Install and Setup ROS2 Humble - Tutorial 1}

To install packages in ROS2, use the following command syntax:
\begin{lstlisting}[style=bashStyle, language=bash]
sudo apt install ros-humble-package_name
\end{lstlisting}

\subsection*{Environment Setup}

To configure the environment, you need to edit the terminal configuration file. Open the file using the GNOME text editor (\code{gedit}):
\begin{lstlisting}[style=bashStyle, language=bash]
gedit ~/.bashrc
\end{lstlisting}
\begin{itemize}
    \item \code{gedit} -- GNOME text editor.
    \item \code{\textasciitilde/.bashrc} -- Terminal configuration file.
\end{itemize}

\vspace{0.3cm}
\noindent Add the following line at the end of the \code{\textasciitilde/.bashrc} file:
\begin{lstlisting}[style=bashStyle, language=bash]
source /opt/ros/humble/setup.bash
\end{lstlisting}

\begin{itemize}
    \item \code{source} -- Applies settings in this file to the currently open terminal.
    \item \code{/opt/ros/humble/setup.bash} -- ROS script which sets up commands, paths, etc.
\end{itemize}

\vspace{0.3cm}
\begin{notebox}[Summary]
The combination of these steps ensures that the \code{ros2} command works always in any opened terminal.
\end{notebox}

\newpage

\section{Start Your First ROS2 Node - Tutorial 2}

\begin{notebox}[Key Concepts]
\textbf{ROS2 Node} -- ROS2 program, which interacts with ROS2 communications and tools. \\
\textbf{ROS2 Package} -- ROS2 Nodes organized in packages.
\end{notebox}

\subsection*{Writer Listener Demo}

To start the demo, use the \code{ros2 run} command followed by the package name and node name:
\begin{lstlisting}[style=bashStyle, language=bash]
ros2 run demo_nodes_cpp talker
\end{lstlisting}
\begin{itemize}
    \item \code{demo\_nodes\_cpp} -- Package name.
    \item \code{talker} -- Node name.
\end{itemize}

\vspace{0.3cm}
\noindent To start the listener node, open a new terminal and run:
\begin{lstlisting}[style=bashStyle, language=bash]
ros2 run demo_nodes_cpp listener
\end{lstlisting}

\subsection*{Turtlesim Demo}

To run the turtlesim demo, start the main node and the teleoperation node:
\begin{lstlisting}[style=bashStyle, language=bash]
ros2 run turtlesim turtlesim_node
ros2 run turtlesim turtle_teleop_key
\end{lstlisting}

\subsection*{ROS2 Tools}

ROS2 tool, shows running Nodes:
\begin{lstlisting}[style=bashStyle, language=bash]
rqt_graph
\end{lstlisting}

\newpage

\section{Create and Set Up a ROS2 Workspace - Tutorial 3}

\textbf{Colcon} -- build tool for ROS2, builds code and Nodes.

\vspace{0.3cm}
\noindent Update the list of software available in external libraries:
\begin{lstlisting}[style=bashStyle, language=bash]
sudo apt update
\end{lstlisting}

\noindent Install the colcon tool:
\begin{lstlisting}[style=bashStyle, language=bash]
sudo apt install python3-colcon-common-extensions
\end{lstlisting}

\subsection*{Colcon Autocompletion}

Add the following line at the end of the \code{\textasciitilde/.bashrc} file to enable autocompletion for colcon arguments:
\begin{lstlisting}[style=bashStyle, language=bash]
source /usr/share/colcon_argcomplete/hook/colcon-argcomplete.bash
\end{lstlisting}

\subsection*{Workspace Setup}

To set up a new workspace, create the workspace folder and the source folder:
\begin{lstlisting}[style=bashStyle, language=bash]
mkdir ros2_ws
mkdir ros2_ws/src
\end{lstlisting}

\vspace{0.3cm}
\noindent Inside the \code{ros2\_ws} folder, run the command to build the ROS2 workspace in the current folder:
\begin{lstlisting}[style=bashStyle, language=bash]
colcon build
\end{lstlisting}

\subsection*{Sourcing the Workspace}

Add the following line at the end of the \code{\textasciitilde/.bashrc} file to allow using custom ROS2 Nodes in any terminal created in this workspace:
\begin{lstlisting}[style=bashStyle, language=bash]
source ~/workspace_name/install/setup.bash
\end{lstlisting}

\newpage

\section{Create a ROS2 Python Package - Tutorial 4}

\vspace{0.3cm}
\noindent \textbf{In \code{src} folder:}

\vspace{0.1cm}
\noindent Create a ROS2 Package:
\begin{lstlisting}[style=bashStyle, language=bash]
ros2 pkg create package-name --build-type
\end{lstlisting}
\noindent Specify if it's a C++/Python package using the build type:
\begin{itemize}
    \item \code{ament\_cmake} -- for C++
    \item \code{ament\_python} -- for Python
\end{itemize}

\vspace{0.3cm}
\noindent Add rcl (ROS Client Library) API:
\begin{lstlisting}[style=bashStyle, language=bash]
--dependencies
\end{lstlisting}
\begin{itemize}
    \item \code{rclcpp} -- for C++
    \item \code{rclpy} -- for Python
\end{itemize}

\vspace{0.3cm}
\begin{notebox}
In every package, there is a \code{package.xml} file containing the name, dependencies, build-type, etc., and a similar \code{setup.py} file.
\end{notebox}

\vspace{0.3cm}
\noindent \textbf{In workspace folder:}

\vspace{0.1cm}
\noindent Build the package, needed before starting nodes:
\begin{lstlisting}[style=bashStyle, language=bash]
colcon build
\end{lstlisting}

\newpage

\section{Create a ROS2 Node with Python and OOP - Tutorial 5}

\vspace{0.3cm}
\noindent \textbf{In \code{package\_name/package\_name}:}

\vspace{0.1cm}
\noindent Create Node:
\begin{lstlisting}[style=bashStyle, language=bash]
touch node_name.py
\end{lstlisting}

\noindent Make Node executable:
\begin{lstlisting}[style=bashStyle, language=bash]
chmod +x node_name.py
\end{lstlisting}

\vspace{0.3cm}
\noindent \textbf{Python Node Implementation:}
\begin{lstlisting}[style=pythonStyle, language=Python]
#!/usr/bin/env python3
# Tell interpreter to use Python3

import rclpy # Import Python library for ROS2
from rclpy.node import Node # Created Nodes inherit from built-in Node class

class MyNode(Node): # MyNode inherits from Node
    def __init__(self):
        super().__init__("first_node") # names Node
        self.counter_ = 0
        self.create_timer(1.0, self.timer_callback)

    def timer_callback(self):
        self.get_logger().info("Hello " + str(self.counter_))
        self.counter_ += 1

def main(args=None):
    rclpy.init(args=args) # initializes ROS2 communications
    node = MyNode()
    rclpy.spin(node) # node will continue to run until killed
    rclpy.shutdown() # terminates ROS2 communications

if __name__ == '__main__': # import protection
    main()
\end{lstlisting}

\vspace{0.3cm}
\noindent \textbf{In \code{setup.py}:}

\vspace{0.1cm}
\noindent Add to entry points:
\begin{lstlisting}[style=pythonStyle, language=Python]
entry_points={
    'console_scripts': [
        'test_node = package_name.file_name:main'
    ],
},
\end{lstlisting}
\begin{itemize}
    \item \code{test\_node} -- executable name.
    \item \code{main} -- function name.
\end{itemize}

\vspace{0.3cm}
\noindent Run updated version of Node:
\begin{lstlisting}[style=bashStyle, language=bash]
colcon build --symlink-install
source ~/.bashrc
ros2 run package_name node_name
\end{lstlisting}
\begin{itemize}
    \item \code{--symlink-install} -- run updated version without build.
\end{itemize}

\vspace{0.3cm}
\noindent List running Nodes:
\begin{lstlisting}[style=bashStyle, language=bash]
ros2 node list
\end{lstlisting}

\noindent Show info about a Node:
\begin{lstlisting}[style=bashStyle, language=bash]
ros2 node info /node_name
\end{lstlisting}

\newpage

\section{What is a ROS2 Topic? - Tutorial 6}

\begin{warnbox}[Communication Flow]
\code{/talker} (Publisher Node) $\rightarrow$ \code{/chatter} (Topic) $\rightarrow$ \code{/listener} (Subscriber Node)
\end{warnbox}

\vspace{0.3cm}
\noindent List running Topics:
\begin{lstlisting}[style=bashStyle, language=bash]
ros2 topic list
\end{lstlisting}

\vspace{0.3cm}
\noindent Show info about a Topic:
\begin{lstlisting}[style=bashStyle, language=bash]
ros2 topic info /topic_name
\end{lstlisting}
\code{Type: std\_msgs/msg/String}\\
\code{Publisher count: 1}\\
\code{Subscription count: 1}

\vspace{0.3cm}
\noindent Show details of a message interface:
\begin{lstlisting}[style=bashStyle, language=bash]
ros2 interface show std_msgs/msg/String
\end{lstlisting}
\code{string data} -- sent message contains data of 'string' type.

\vspace{0.3cm}
\noindent Show data received by Topic:
\begin{lstlisting}[style=bashStyle, language=bash]
ros2 topic echo /chatter
\end{lstlisting}

\vspace{0.5cm}
\begin{notebox}
There can be multiple Nodes publishing to the same Topic and multiple Nodes subscribing to the same topic.
\end{notebox}

\newpage

\section{Write a ROS2 Publisher with Python - Tutorial 7}

\vspace{0.3cm}
\noindent Import message type:
\begin{lstlisting}[style=pythonStyle, language=Python]
from geometry_msgs.msg import Twist
\end{lstlisting}

\vspace{0.3cm}
\noindent In \code{package.xml} add imported module:
\begin{lstlisting}[style=xmlStyle]
<depend>geometry_msgs</depend>
\end{lstlisting}

\vspace{0.3cm}
\noindent \textbf{In constructor:}
\begin{lstlisting}[style=pythonStyle, language=Python]
self.cmd_vel_pub_ = self.create_publisher(Twist, "/turtle1/cmd_vel", 10)
\end{lstlisting}
\begin{itemize}
    \item \code{self.cmd\_vel\_pub\_} -- Create publisher.
    \item \code{Twist} -- Message type.
    \item \code{"/turtle1/cmd\_vel"} -- Topic name.
    \item \code{10} -- queue size.
\end{itemize}

\vspace{0.3cm}
\noindent \code{geometry\_msgs/msg/Twist} consists of:
\begin{itemize}
    \item \code{linear} (position)
    \item \code{angular} (rotation)
\end{itemize}

\vspace{0.3cm}
\noindent \textbf{Publishing function:}
\begin{lstlisting}[style=pythonStyle, language=Python]
def send_velocity_command(self):
    msg = Twist()
    msg.linear.x = 2.0
    msg.angular.z = 1.0
    self.cmd_vel_pub_.publish(msg) # send msg via publisher
\end{lstlisting}

\newpage

\section{Write a ROS2 Subscriber with Python - Tutorial 8}

\vspace{0.3cm}
\noindent Import message type:
\begin{lstlisting}[style=pythonStyle, language=Python]
from turtlesim.msg import Pose
\end{lstlisting}

\vspace{0.3cm}
\noindent In \code{package.xml} add imported module:
\begin{lstlisting}[style=xmlStyle]
<depend>turtlesim</depend>
\end{lstlisting}

\vspace{0.3cm}
\noindent \textbf{In constructor:}
\begin{lstlisting}[style=pythonStyle, language=Python]
self.pose_publisher_ = self.create_subscription(Pose, "/turtle1/pose", self.pose_callback, 10)
\end{lstlisting}
\begin{itemize}
    \item \code{self.create\_subscription} -- Create subscriber.
    \item \code{Pose} -- Message type.
    \item \code{"/turtle1/pose"} -- Topic name.
    \item \code{self.pose\_callback} -- Callback name.
    \item \code{10} -- Queue size.
\end{itemize}

\vspace{0.3cm}
\begin{notebox}
Callback will be called everytime when you receive a message on that Topic.
\end{notebox}

\vspace{0.3cm}
\noindent \textbf{Callback function receiving messages:}
\begin{lstlisting}[style=pythonStyle, language=Python]
def pose_callback(self, msg: Pose):
    self.get_logger().info("(" + str(msg.x) + ", " + str(msg.y) + ")")
\end{lstlisting}

\newpage

\section{Create a Closed Loop System with a Publisher and a Subscriber - Tutorial 9}

\vspace{0.3cm}
\noindent Combine previous two tutorials -- Create a Node that both subscribes and publishes appropriate command when received a message.

\newpage

\section{What is a ROS2 Service? - Tutorial 10}

\vspace{0.3cm}
\noindent Services are interactions that act like a Client/Server. Client sends a request to the Server and Server responds to the Client. Topics don't give answers.

\vspace{0.3cm}
\noindent Services cannot be seen in \code{rqt\_graph}. List running services:
\begin{lstlisting}[style=bashStyle, language=bash]
ros2 service list
\end{lstlisting}

\vspace{0.3cm}
\noindent See Service type:
\begin{lstlisting}[style=bashStyle, language=bash]
ros2 service type /service_name
\end{lstlisting}
Example: \code{example\_interfaces/srv/AddTwoInts}

\vspace{0.3cm}
\noindent See data inside Service message:
\begin{lstlisting}[style=bashStyle, language=bash]
ros2 interface show example_interfaces/srv/AddTwoInts
\end{lstlisting}
\code{int64 a} -- request \\
\code{int64 b} -- request \\
\code{---} \\
\code{int64 sum} -- response

\vspace{0.3cm}
\noindent Make a request:
\begin{lstlisting}[style=bashStyle, language=bash]
ros2 service call /service_name service/type "{'a': 2, 'b': 5}"
\end{lstlisting}
\begin{itemize}
    \item Response on client from server: \code{(...) sum = 7}
    \item Info about request on server: \code{(...) a: 2 b: 5}
\end{itemize}

\vspace{0.5cm}
\begin{warnbox}[Server/Client Architecture]
There can only be one server, but many clients connecting to the server.
\end{warnbox}

\vspace{0.3cm}
\noindent \textbf{Requests} are typically used for:
\begin{itemize}
    \item Computation
    \item Change of settings
\end{itemize}

\vspace{0.3cm}
\begin{notebox}[Topics vs Services]
Topics are used to send straight data, without expecting an answer.
\end{notebox}

\newpage

\section{Write a ROS2 Service Client with Python - Tutorial 11}

\vspace{0.3cm}
\noindent Import module:
\begin{lstlisting}[style=pythonStyle, language=Python]
from functools import partial
\end{lstlisting}

\vspace{0.3cm}
\noindent In \code{package.xml} add imported module:
\begin{lstlisting}[style=xmlStyle]
<depend>functools</depend>
\end{lstlisting}

\vspace{0.3cm}
\noindent \textbf{Create a Service-sending function:}
\begin{lstlisting}[style=pythonStyle, language=Python]
def call_set_pen_service(self, argument):
    client = self.create_client(ServiceType, "/service_name")
    
    # Wait for service to be available
    while not client.wait_for_service(1.0):
        self.get_logger().warn("Waiting for service")
        
    request = ServiceType.Request()
    request.argument = argument
    
    # Request without waiting for response
    future = client.call_async(request)
    future.add_done_callback(partial(self.callback))
\end{lstlisting}

\vspace{0.3cm}
\begin{notebox}[Using partial]
\code{partial} is used when, for example, a callback needs more arguments than ROS can give at the moment. \code{partial} allows adding missing arguments in the meantime.
\end{notebox}

\vspace{0.3cm}
\noindent \textbf{Callback function:}
\begin{lstlisting}[style=pythonStyle, language=Python]
def callback(self, future):
    try:
        response = future.result()
    except Exception as e:
        self.get_logger().error("Response error")
\end{lstlisting}

\vspace{0.3cm}
\begin{warnbox}[Topics vs Services Usage]
You can publish on a Topic very frequently. A Service is made to be called only when needed. Otherwise, it can really slow down the application.
\end{warnbox}

\vspace{0.3cm}
\noindent Get frequency of a Topic:
\begin{lstlisting}[style=bashStyle, language=bash]
ros2 hz /topic_name
\end{lstlisting}

\end{document}